% \documentclass[manuscript,review,anonymous]{acmart}
\documentclass[sigconf]{acmart}

\AtBeginDocument{%
  \providecommand\BibTeX{{%
    \normalfont B\kern-0.5em{\scshape i\kern-0.25em b}\kern-0.8em\TeX}}}

\setcopyright{rightsretained}

\copyrightyear{2026}
\acmYear{2026}
\setcopyright{cc}
\setcctype{by}
\acmConference[CHI '26]{Proceedings of the 2026 CHI Conference on Human Factors in Computing Systems}{April 13--17, 2026}{Barcelona, Spain}
\acmBooktitle{Proceedings of the 2026 CHI Conference on Human Factors in Computing Systems (CHI '26), April 13--17, 2026, Barcelona, Spain}
\acmPrice{}
\acmDOI{10.1145/3772318.3790375}
\acmISBN{979-8-4007-2278-3/2026/04}

\newcommand\subjectcount{246}


\newcommand{\todo}[1]{{\color{blue} #1}}
\newcommand{\del}[1]{\textcolor{red}{#1}}

% \usepackage{xeCJK}
\usepackage{multirow}
\usepackage{subcaption}
\usepackage{caption}
\usepackage[normalem]{ulem}
\usepackage{siunitx}
\usepackage{booktabs}
\usepackage{array}
\usepackage{tabularx}
\usepackage{xcolor}
\usepackage{multicol}
\usepackage{siunitx}
\usepackage{pifont}
\usepackage[utf8]{inputenc}

% Shared box style 
\newcommand{\SharedGrayBox}[2]{%
  \par\vspace{4pt}
  \noindent
  \hspace*{-6pt}% shift left (adjust value as needed)
  \begingroup
  \setlength{\fboxsep}{5pt}      % inner padding
  \setlength{\fboxrule}{1pt}      % border thickness
  \fcolorbox{gray!75!black}{gray!5}{%
    \begin{minipage}{0.98\linewidth}
      \textbf{#1}\par
      \vspace{6pt}
      {\ttfamily\footnotesize
      #2\par
      }
    \end{minipage}
  }%
  \endgroup
  \par\vspace{4pt}
}


\begin{document}

\title[Towards AI as Colleagues: Multi-Agent System Improves Structured Ideation Processes]{Towards AI as Colleagues:  \\ Multi-Agent System Improves Structured Ideation Processes}

\author{Kexin Quan}
\affiliation{%
  \institution{Information Sciences, \\University of Illinois Urbana-Champaign}
  \city{Champaign}
  \state{Illinois}
  \country{USA}}
\email{kq4@illinois.edu}

\author{Dina Albassam}
\authornote{Equal contribution.}
\affiliation{%
  \institution{Computer Science, \\University of Illinois Urbana-Champaign}
  \city{Champaign}
  \state{Illinois}
  \country{USA}}
\email{dinasa2@illinois.edu}

\author{Mengke Wu}
\authornotemark[1]
\affiliation{%
  \institution{Information Sciences, \\University of Illinois Urbana-Champaign}
  \city{Champaign}
  \state{Illinois}
  \country{USA}}
\email{mengkew2@illinois.edu}


\author{Zijian Ding}
\affiliation{%
  \institution{College of Information, \\University of Maryland}
  \city{College Park}
  \state{Maryland}
  \country{USA}}
\email{ding@umd.edu}

\author{Jessie Chin}
\affiliation{%
  \institution{Information Sciences, \\University of Illinois Urbana-Champaign}
  \city{Champaign}
  \state{Illinois}
  \country{USA}}
\email{chin5@illinois.edu}

\renewcommand{\shortauthors}{Quan, et al.}


\begin{abstract}
Most AI systems today are designed to manage tasks and execute predefined steps. This makes them effective for process coordination but limited in their ability to engage in joint problem-solving with humans or contribute new ideas. We introduce MultiColleagues, a multi-agent conversational system that shows how AI agents can act as colleagues by conversing with each other, sharing new ideas, and actively involving users in collaborative ideation processes. In a within-subjects study with 20 participants, we compared MultiColleagues to a single-agent baseline. Results show that MultiColleagues fostered stronger perceived social presence, and participants rated their outcomes as higher in quality and novelty, with more elaboration during ideation. These findings demonstrate the potential of AI agents to move beyond process partners toward colleagues that share intent, strengthen group dynamics, and collaborate with humans to advance ideas.
\end{abstract}

%%
%% The code below is generated by the tool at http://dl.acm.org/ccs.cfm.
%% Please copy and paste the code instead of the example below.
%%
\begin{CCSXML}
<ccs2012>
   <concept>
       <concept_id>10003120.10003121.10011748</concept_id>
       <concept_desc>Human-centered computing~Empirical studies in HCI</concept_desc>
       <concept_significance>500</concept_significance>
       </concept>
   <concept>
       <concept_id>10003120.10003121.10003124.10011751</concept_id>
       <concept_desc>Human-centered computing~Collaborative interaction</concept_desc>
       <concept_significance>300</concept_significance>
       </concept>
   <concept>
       <concept_id>10003120.10003130.10011762</concept_id>
       <concept_desc>Human-centered computing~Empirical studies in collaborative and social computing</concept_desc>
       <concept_significance>500</concept_significance>
       </concept>
   <concept>
       <concept_id>10003120.10003123</concept_id>
       <concept_desc>Human-centered computing~Interaction design</concept_desc>
       <concept_significance>500</concept_significance>
       </concept>
 </ccs2012>
\end{CCSXML}

\ccsdesc[500]{Human-centered computing~Empirical studies in HCI}
\ccsdesc[300]{Human-centered computing~Collaborative interaction}
\ccsdesc[500]{Human-centered computing~Empirical studies in collaborative and social computing}
\ccsdesc[500]{Human-centered computing~Interaction design}

\keywords{AI as Colleagues, Multi-Agent Systems, Large Language Models, Interaction Design, Collaborative Ideation}


\begin{teaserfigure}
\centering
\includegraphics[width=0.79\textwidth]{figures/teaser.jpg}
\caption[Core scenario of MultiColleagues.]{\textbf{MultiColleagues: An Illustrative User Scenario of AI-Supported Ideation.}
P10 introduces a karaoke-in-autonomous-vehicle challenge (Step 1), selects three AI colleagues (Step 2), then moves from \textit{Explore} (divergent ideas) to \textit{Focus} (convergence) on \textit{Trustworthy mood-aware Karaoke} with privacy, simplicity, and local data processing (Step 3).}
\Description{Participant 10 scenario with MultiColleagues. The figure depicts Participant 10 identifying a challenge, selecting role-differentiated colleagues, and joining the discussion. Through explore and focus modes, divergent and convergent thinking are applied to generate and refine ideas.}
\end{teaserfigure}


\maketitle

\def \RQO {\textbf{RQ1}: ?}

\def \RQT {\textbf{RQ2}: ?}

\input{sections/1-introduction}

\input{sections/2-relatedWork}

\input{sections/3-studyDesign}

\input{sections/4-results}

\input{sections/5-discussion_conclusion}

\end{document}
\endinput

